\section{模型的评价、改进与推广}

\subsection{模型的创新之处}

本文的第一处创新是界面扩散系数的取法。有限体积法在界面处如何取扩散系数通常被视为实现细节，本题却因扩散系数跨越四个数量级而使这一细节变成决定性的：调和平均在相邻节点含水率悬殊时被小值一侧完全支配，实算给出的停时随网格加密单调漂移而不收敛（$N=20/40/80$ 依次为 $549.918$、$223.428$、$87.440$~h，每加倍减半）。本文改按式~\eqref{eq:iface} 取扩散系数在两节点之间的积分平均，即以 Kirchhoff 势的差商定义界面通量，网格加密下停时立刻收敛（$57.562$、$57.532$、$57.524$~h）。这一取法在物性真正分区跳变时自动退化为调和平均，因此不是以牺牲适用性换取收敛，而是调和平均与算术平均的共同推广。

第二处创新是收缩域上的物质坐标处理。式~\eqref{eq:mapping} 以当前半径归一，使表面永远落在坐标端点 $s=1$，边界条件的施加位置不随时间改变；更关键的是，链式法则产生的坐标变换项与表观对流项精确相消，最终方程中不出现半径的导数，因而不必对实测半径序列作数值微分——附件 2 的半小时采样若直接差分，会引入阶梯状脉冲。收缩效应完全通过 $1/R^2(t)$、边界系数中的 $1/R(t)$ 与控制体的物理体积三处进入方程，物理机制反而更清楚：半径由 $2$~cm 减到 $1.2$~cm 使扩散项放大 $2.78$ 倍，这是收缩加速干燥的主要来源。

第三处创新是把守恒校核用作实现的诊断手段而非事后装点。式~\eqref{eq:massres} 在本文的离散下是代数恒等式，残差只反映舍入，因此它对编码疏漏极其灵敏。实验室坐标对照支的停时与主模型只差 $2.93\%$，凭数值本身无法判断对错，而其守恒残差高达 $30.3\%$，据此才定位到该支把边界运动计入了两次。同一思路也用在收缩方式的判定上：由干物质守恒式~\eqref{eq:drymass} 反推的干物质密度与附录 4 密度公式的两个自然参照相比，径向仿射收缩相容到十几个百分点以内而各向同性收缩偏离近一倍，这一判断只用到附件 2 的几何序列与附录 4 的密度公式两项互不相干的题给信息。

第四处是效应分解的做法。问题四与问题三同时差在物性组与几何两处，直接比较停时只能得到 $11.11\%$ 的净差，看不出内部结构。本文追加一个只换物性而保持半径不变的隔离算例，把净效应拆成物性使干燥慢一倍以上（$+125.8\%$）与收缩使干燥快约六成（$-60.6\%$）两项方向相反的贡献，从而说明净差的符号不能由任何一侧的定性推理单独断定。

\subsection{模型的不足}

最需要说明的一项是热方程未计蒸发潜热。题目给出的四组物性附录都没有潜热数据，本文按假设（3）不设潜热汇，但这一选择的影响是全部对照算例中最大的：加入 $2.4\times10^6$~J/kg 的潜热汇后停时变为 $62.9380$~h，比基准慢 $9.40\%$。本文不宣称这一项可以忽略，只能指出若题目提供潜热数据，停时应向偏慢一侧修正约一成，而结论的量级不变。

其次是一维径向的降维。药材长径比为 $6.25$，两端面积占总表面的 $7.4\%$，端面散失被忽略会使算得的停时偏大。本文未实算二维轴对称模型，因此只能从比例上论证这一项的量级，不能给出精确修正；可以确定的方向是端部区域会先达标，而题目所问的“到药材中心的距离”正对应本文求解的中截面。

第三是收缩域上能量的处理。水分方程在物质坐标下严格守恒，因为其测度计的是不随收缩改变的干物质份额；热容项的物理测度含有 $R^2(t)$，本文的离散未追加 $\mathrm{d}(R^2)/\mathrm{d}t$ 一项，故收缩域上的能量并非严格守恒。本文未直接估计该项的量值；可以说的只是温度耦合本身很弱——把温度场整段冻结为空气温度只使停时变化 $0.15\%$——因此这一疏漏难以经由温度耦合显著改变停时。

第四是空气条件的时段延拓。附件 1 覆盖到 $14400$~s，而干燥持续二十万秒，其余时间的空气条件按实测尾段的平台值取定。判断依据是尾段 $61$ 点的逐点增量已落在零线附近，且空气水分浓度 $\pm10\%$ 的扰动只带来 $0.2\%$ 的效应，因此这一处理不是结论的薄弱环节；但严格地说，平台之外若发生工况变动，本文无从反映。

最后，表面传质驱动力的参考基存在两种解释，本文按量纲一致性取药材侧与空气侧含水率的直接差值。在问题三的算例上改用另一种解释，停时相差 $1.21\%$，达标时的表面含水率由 $0.0526$ 压到 $0.0026$~kg/kg。轴线达标时刻几乎不受影响，但表面的绝对水平确实依赖于这一选择。

\subsection{模型的改进与推广}

在题给数据之内，最直接的改进是把水分迁移由单一等效 Fick 扩散拆分为自由水与结合水两种机制，或者按 Luikov 方程组同时考虑温度梯度驱动的水分迁移；这两条路线都需要本题未提供的参数，因此只能作为方向指出。若能获得潜热与表面蒸发速率的数据，热方程加入潜热汇是收益最明确的一步改进，量级已由前述对照算例给出。若要消除端面效应，把当前离散扩展为二维轴对称 $(r,z)$ 网格并不需要改变任何物理设定，只是计算量按节点数平方增长。

本文的框架不依赖药材这一具体对象。方程组、界面取法、停时判据与移动边界处理只用到三件事：一维几何、随解变化的物性、第三类边界，因此可以直接搬到其他以内部扩散为限制环节的干燥与浸取过程，例如谷物与果蔬的热风干燥、多孔颗粒的溶剂脱除；热风干燥中失水进程与成品质地的关联也说明，能给出内部分布而非仅给出平均值的模型对工艺评价更有用\upcite{ashtiani2017hotair}。判据一侧的推广更宽：本文把“各处合格”实现为对径向最大值的首次穿越判定，这一写法对最湿位置的所在不作先验假设，因此在存在内部极值或多个失水面的工况下同样成立。移动边界的物质坐标变换只要求区域的变形是仿射的，对收缩率的具体形式没有限制，实测序列、经验公式或与含水率联动的收缩模型都可以代入同一套离散。

就工程使用而言，本文的灵敏度排序给出了一条清晰的取舍：预测精度主要取决于扩散系数与潜热数据的可靠程度，表面传热系数的测量精度几乎不影响结果；而合格标准本身的效应超过所有物理参数，且收紧标准的代价明显大于放宽标准的收益——把阈值收紧一成要多烘三成时间。这一条对烘房的实际操作比任何单个参数的精确值都更有价值。

% 数据保真核对：摘要、正文、附录表格中的结果/比较/误差/最优值已逐项对照
% figures/*.json 与 code/params.py 实算产物；共核 459 个高精度数字 + 约 420 个表格单元。
% 已修正 6 处：11_compare 倍数配对、6_problem1 的 D(C0) 与 alpha 舍入、11_compare 的 Bi_m 区间、
% 5_framework 的平台步长收敛值与分段步长、10_credibility 的 R^2 舍入。
% DATA_CHECK_PASSED
